Die Scharniere sind die Drehachse. Drückst du senkrecht zur Tür, ist der Hebelarm der Abstand $d$ von den Scharnieren, und es gilt $M = F\cdot d$, also $F = M/d$ (Längen in Metern).

\begin{abcliste}
\abc
An der Klinke:
\begin{align}
 F_1 &= \FaF \\
   &= \frac{\M}{\da} \\
   &= \Fa \approx \result{\FaP{0}{2}}
\end{align}

\abc
Nahe beim Scharnier:
\begin{align}
 F_2 &= \FbF \\
   &= \frac{\M}{\db} \\
   &= \Fb \approx \result{\FbP{0}{2}}
\end{align}
Ein kürzerer Hebelarm braucht eine grössere Kraft: Darum sind Klinken weit weg von den Scharnieren.
\end{abcliste}
